\mode*
\section{Conclusions}
\label{conclusions}

We set out to catalogue the misconceptions that students bring to introductory
programming and to analyse how to teach against them. The catalogue
(\cref{misconceptions}) combines a manual literature review with a systematic,
tool-assisted search, and is complemented by misconceptions we observed in our
own teaching. For each misconception we identified the object of learning and
its critical aspects, and derived a pattern of variation --- contrast,
generalisation and, where several critical aspects interact, fusion --- that
makes the critical aspects discernible.

The payoff of the variation-theory analysis is that it prescribes the
\emph{content} of teaching at the level of examples: which programs to
juxtapose, what to vary between them, and what to hold invariant, so that
precisely the misconceived aspect comes into focus. This complements
activity-level teaching designs, as discussed in \cref{relatedwork}, which can
serve as carriers for the derived patterns. Variation theory moreover predicts
that students taught through patterns of variation become better at generating
such patterns themselves, an effect that surfaces as delayed, generative
learning \autocite[pp.~227--228]{NCOL}; since debugging draws on the same
ability --- varying one aspect while holding others invariant to isolate an
unknown --- we hypothesise that teaching with these patterns also improves
students' debugging.

Our contribution is analytical, and its limitations mark out the work that
remains. The derived patterns have not yet been evaluated in the classroom.
The natural design for doing so is a \emph{learning study}
\autocite[Ch.~7]{NCOL} for each topic --- or even for each pattern --- with
pre- and post-tests that test what is to be taught; we plan such a study in
the coming iteration of our introductory programming course, and we derive
the diagnostic instruments for it from the catalogue in
\cref{app:diagnostics}: one item per critical aspect, each distractor
encoding a catalogued misconception. Instrumenting
the students' own work, to log how they use variation during their studies,
would moreover let us test the debugging hypothesis above directly.
The literature coverage has recorded
limits: one major index was unavailable during the systematic search, the
citation snowball ran for one generation only (\cref{app:litlimitations}),
and several papers kept by the screenings remain to be integrated into the
catalogue. Some
references are, transparently, verified at the abstract level pending
institutional access to the full texts. Finally, the misconceptions we
observed in our own teaching should be read as hypotheses for teaching
designs rather than established findings, until they are corroborated the
way the rest of the catalogue is.

\mode<presentation>{%
\begin{frame}
  \begin{block}{Conclusions}
    \begin{itemize}
      \item Catalogue of CS1 misconceptions: literature (manual + systematic)
        and our own teaching.
      \item Each analysed through variation theory: critical aspects, then
        contrast--generalisation--fusion.
      \item The analysis prescribes example \emph{content}; activity designs
        (PRIMM, peer instruction, \ldots) can carry it.
      \item Hypothesis: generative learning transfers to debugging.
      \item Next: integrate remaining papers; a learning study per topic
        (pre-/post-tests); log students' use of variation.
    \end{itemize}
  \end{block}
\end{frame}%
}
